\begin{chapter}{Equidade das decisões}\label{cap:estudo5}

Este capítulo corresponde ao quinto objetivo específico: avaliar o desempenho das decisões por subgrupo, quanto à paridade de sensibilidade e de valor preditivo positivo, explicitando a incompatibilidade geral desses critérios sob prevalências desiguais, distinguindo subgrupos demográficos de subgrupos operacionais, como o quartil de concordância do nome e a origem do registro, e registrando a ausência da variável raça/cor nas bases utilizadas. Apresentam-se o método, os resultados e a discussão desse estudo.

\section{Método}\label{sec:e5-metodo}

Para investigar se o desempenho do classificador varia de forma sistematicamente desigual entre subgrupos demográficos e epidemiológicos representados na base de pares candidatos, realizou-se uma análise de equidade inspirada no \textit{framework} de \textit{equalized odds} \cite{Hardt2016equality}. Adotou-se, como critério operacional, a paridade simultânea da taxa de verdadeiros positivos (TVP) e do valor preditivo positivo (VPP), combinação que captura tanto a capacidade de identificar pares verdadeiros quanto a proporção de predições positivas corretas em cada estrato \cite{Barocas2019fairml}.

Cinco dimensões de estratificação foram definidas a partir de atributos disponíveis na base de pares. (a)~Quartis de concordância do campo nome, calculados sobre o sub-escore \texttt{NOME qtd frag iguais}, que representa a proporção de fragmentos idênticos entre os nomes dos dois registros, dividida pela média do total de fragmentos de ambos (cf.\ Seção~\ref{sec:comparador-registros}); essa variável foi adotada como \textit{proxy} da complexidade onomástica do par, uma vez que nomes com poucos fragmentos em comum tendem a ser mais difíceis de classificar e podem revelar viés diferencial do modelo. (b)~Origem do par (controle \textit{versus} TB-origem), distinguindo pares cujo registro SIM tem tuberculose como causa de óbito daqueles em que a TB figura apenas no Sinan-TB. (c)~Faixa etária em cinco categorias (0--17, 18--39, 40--59, 60--79, 80+ anos), correspondentes aos estratos epidemiológicos empregados na análise da carga de TB no Brasil \cite{Brasil2024tb}, que capturam diferenças no perfil de notificação e na completude dos campos de identificação entre crianças, adultos jovens, adultos, idosos e longevos. (d)~Sexo (feminino \textit{versus} masculino). (e)~Município de residência, representado pelo código IBGE do campo \texttt{CODMUNRES}; cabe notar que, na presente base, mais de 99\% dos registros provêm de um único município (Rio de Janeiro), de modo que essa dimensão avalia essencialmente o contraste entre o município-sede e os demais da região metropolitana, limitando a generalização de eventuais diferenças geográficas.

Para cada dimensão, a TVP e o VPP foram estimados separadamente por grupo (e.g., por faixa etária), acompanhados de intervalos de confiança de Wilson. A comparação entre estratos foi conduzida par a par: para cada combinação de dois grupos dentro da mesma dimensão (e.g., faixa 0--17 \textit{versus} 18--39, faixa 0--17 \textit{versus} 40--59, e assim por diante), construiu-se uma tabela de contingência $2 \times 2$ cruzando o grupo de pertencimento com o resultado da classificação (acerto ou erro). Quando todas as frequências esperadas dessa tabela eram iguais ou superiores a cinco, empregou-se o teste qui-quadrado de Pearson; nos demais casos, o teste exato de Fisher \cite{Agresti2007}. Para controlar a inflação da taxa de falsas descobertas decorrente das múltiplas comparações, aplicou-se a correção de Benjamini-Hochberg ao nível $\alpha = 0{,}05$ sobre o conjunto completo de testes \cite{Benjamini1995fdr}. A magnitude da eventual disparidade entre dois grupos foi quantificada por dois indicadores complementares: a razão de chances (\textit{odds ratio}) de classificação correta, que expressa quantas vezes a chance de acerto em um grupo excede a do outro, e a diferença absoluta de taxas (TVP ou VPP entre os dois estratos), que fornece a distância em pontos percentuais entre os desempenhos.

\section{Resultados}\label{sec:e5-resultados}

\subsection{Análise de equidade por subgrupos demográficos e de similaridade}\label{sec:e5-resultados-equidade}

O \textit{framework} de equidade, incluindo o critério adotado, as dimensões de estratificação e os testes estatísticos empregados, é descrito na Seção~\ref{sec:e5-metodo}.

\begin{table}[htbp]
\centering
\caption{Sensibilidade e valor preditivo positivo (VPP), em porcentagem, por subgrupo: XGBoost com limiar de sensibilidade equivalente ao corte 6{,}0 e decis\~ao autom\'atica da cascata v3.1 (vigil\^ancia).}
\label{tab:qw3-group-metrics}
\small
\begin{tabular}{llrrrrrr}
\toprule
 & & & & \multicolumn{2}{c}{XGBoost} & \multicolumn{2}{c}{Cascata v3.1} \\
\cmidrule(lr){5-6}\cmidrule(lr){7-8}
Dimens\~ao & Grupo & $n$ & Vínculos & Sens. & VPP & Sens. & VPP \\
\midrule
\multicolumn{8}{l}{\textit{Subgrupos demogr\'aficos}} \\
Sexo & F & 12.834 & 71 & 94{,}4 & 93{,}1 & 81{,}7 & 87{,}9 \\
Sexo & M & 48.862 & 176 & 90{,}3 & 70{,}7 & 77{,}3 & 84{,}0 \\
\midrule
Faixa et\'aria & 0 a 17 & 516 & 4 & 75{,}0 & 50{,}0 & 50{,}0 & 50{,}0 \\
Faixa et\'aria & 18 a 39 & 15.718 & 87 & 90{,}8 & 73{,}8 & 75{,}9 & 84{,}6 \\
Faixa et\'aria & 40 a 59 & 24.162 & 95 & 92{,}6 & 84{,}6 & 80{,}0 & 87{,}4 \\
Faixa et\'aria & 60 a 79 & 18.746 & 50 & 90{,}0 & 72{,}6 & 78{,}0 & 92{,}9 \\
Faixa et\'aria & 80+ & 2.151 & 10 & 100{,}0 & 58{,}8 & 100{,}0 & 62{,}5 \\
\midrule
\multicolumn{8}{l}{\textit{Subgrupos operacionais}} \\
Nome (quartis) & Q1 & 15.640 & 1 & 100{,}0 & 100{,}0 & 0{,}0 & n.a. \\
Nome (quartis) & Q2 & 15.883 & 2 & 0{,}0 & 0{,}0 & 0{,}0 & 0{,}0 \\
Nome (quartis) & Q3 & 25.637 & 21 & 61{,}9 & 43{,}3 & 47{,}6 & 55{,}6 \\
Nome (quartis) & Q4 & 4.536 & 223 & 95{,}1 & 80{,}3 & 82{,}5 & 88{,}9 \\
\midrule
Encerramento da notifica\c{c}\~ao & Demais situa\c{c}\~oes & 57.555 & 143 & 89{,}5 & 66{,}0 & 74{,}8 & 79{,}9 \\
Encerramento da notifica\c{c}\~ao & Encerrada por \'obito & 4.141 & 104 & 94{,}2 & 95{,}1 & 83{,}7 & 92{,}6 \\
\midrule
Munic\'ipio & Capital (RJ) & 61.125 & 226 & 91{,}6 & 74{,}5 & 77{,}4 & 85{,}0 \\
Munic\'ipio & Demais & 571 & 21 & 90{,}5 & 100{,}0 & 90{,}5 & 86{,}4 \\
\midrule
Passo & 11 & 57.598 & 40 & 75{,}0 & 35{,}3 & 45{,}0 & 50{,}0 \\
Passo & 6 & 2.649 & 80 & 92{,}5 & 84{,}1 & 76{,}2 & 91{,}0 \\
Passo & 9 & 1.449 & 127 & 96{,}1 & 98{,}4 & 90{,}6 & 92{,}0 \\
\midrule
Ano do \'obito & 2019 & 12.397 & 53 & 88{,}7 & 66{,}2 & 66{,}0 & 92{,}1 \\
Ano do \'obito & 2020 & 13.412 & 52 & 90{,}4 & 78{,}3 & 82{,}7 & 78{,}2 \\
Ano do \'obito & 2021 & 19.923 & 63 & 95{,}2 & 73{,}2 & 81{,}0 & 81{,}0 \\
Ano do \'obito & 2022 & 15.964 & 79 & 91{,}1 & 85{,}7 & 82{,}3 & 90{,}3 \\
\midrule
Completude & Completo & 59.823 & 226 & 91{,}2 & 75{,}2 & 77{,}9 & 84{,}6 \\
Completude & Incompleto & 1.873 & 21 & 95{,}2 & 87{,}0 & 85{,}7 & 90{,}0 \\
\bottomrule
\end{tabular}
\end{table}


A Tabela~\ref{tab:qw3-group-metrics} apresenta as métricas por subgrupo.
A TVP manteve-se estável na maioria das comparações, com rejeição da hipótese nula em apenas um par: faixa etária 18--39 \textit{versus} 40--59 ($p_{\text{FDR}} = 0{,}039$; TVP de $0{,}851$ e $0{,}958$, respectivamente), configurando disparidade limítrofe cuja relevância prática depende do tamanho do denominador (87 e 95 positivos verdadeiros).
Essa relativa homogeneidade da TVP indica que a capacidade do classificador de detectar pares verdadeiros não varia de forma sistemática entre os subgrupos avaliados.

O cenário do VPP revelou-se distinto.
Onze dos doze testes com rejeição da hipótese nula ($p_{\text{FDR}} < 0{,}05$) recaíram sobre o VPP, distribuídos em quatro das cinco dimensões avaliadas (Tabela~\ref{tab:qw3-fairness-rejections}).
Na dimensão de concordância do nome, o VPP do quarto quartil (Q4: $0{,}134$) superou o dos quartis inferiores (Q1--Q3: $< 0{,}004$) por ampla margem, refletindo o fato de que pares com alta similaridade de nome concentram a quase totalidade dos verdadeiros positivos.
Na dimensão origem do registro, registros de origem tuberculosa apresentaram VPP de $0{,}170$, contra $0{,}019$ do grupo controle (razão de chances $= 0{,}094$; $p_{\text{FDR}} < 10^{-88}$), evidenciando que o classificador opera com precisão substancialmente maior no estrato mais relevante para a vigilância.
Na dimensão sexo, o VPP feminino ($0{,}049$) superou o masculino ($0{,}026$) com significância estatística ($p_{\text{FDR}} < 10^{-4}$; RC~$= 1{,}92$), embora a magnitude absoluta da diferença seja modesta.
Na faixa etária, o grupo 0--17 anos exibiu VPP de $0{,}800$, contrastando com valores entre $0{,}017$ e $0{,}042$ nas demais faixas; todavia, o denominador reduzido desse grupo (cinco predições positivas) limita a generalização do achado.

\begin{table}[htbp]
\centering
\caption{Rejei\c{c}\~oes da hip\'otese de paridade (Benjamini-Hochberg, $\alpha = 0{,}05$, 58 testes): XGBoost com limiar de sensibilidade equivalente ao corte 6{,}0. $\Delta$ em pontos percentuais (grupo A menos grupo B).}
\label{tab:qw3-fairness-rejections}
\footnotesize\setlength{\tabcolsep}{3pt}
\begin{tabular}{lllllrrr}
\toprule
Dimens\~ao & Tipo & M\'etrica & Grupo A & Grupo B & $p_{\mathrm{FDR}}$ & $\Delta$ & OR \\
\midrule
Sexo & Demogr. & VPP & F & M & $<$0{,}001 & $+$22{,}4 & 5{,}56 \\
Ano do \'obito & Operac. & VPP & 2019 & 2022 & 0{,}024 & $-$19{,}5 & 0{,}33 \\
Encerramento da notifica\c{c}\~ao & Operac. & VPP & Demais situa\c{c}\~oes & Encerrada por \'obito & $<$0{,}001 & $-$29{,}2 & 0{,}10 \\
Munic\'ipio & Operac. & VPP & Capital (RJ) & Demais & 0{,}046 & $-$25{,}5 & 0{,}07 \\
Nome (quartis) & Operac. & VPP & Q3 & Q4 & $<$0{,}001 & $-$37{,}0 & 0{,}19 \\
Nome (quartis) & Operac. & Sens. & Q3 & Q4 & $<$0{,}001 & $-$33{,}2 & 0{,}08 \\
Nome (quartis) & Operac. & Sens. & Q2 & Q4 & 0{,}020 & $-$95{,}1 & 0{,}01 \\
Passo & Operac. & VPP & 11 & 9 & $<$0{,}001 & $-$63{,}1 & 0{,}01 \\
Passo & Operac. & VPP & 11 & 6 & $<$0{,}001 & $-$48{,}8 & 0{,}10 \\
Passo & Operac. & VPP & 6 & 9 & $<$0{,}001 & $-$14{,}3 & 0{,}09 \\
Passo & Operac. & Sens. & 11 & 9 & 0{,}002 & $-$21{,}1 & 0{,}12 \\
Passo & Operac. & Sens. & 11 & 6 & 0{,}041 & $-$17{,}5 & 0{,}24 \\
\bottomrule
\end{tabular}
\end{table}


A interpretação conjunta desses resultados aponta para um padrão coerente: as disparidades de VPP decorrem predominantemente da distribuição desigual de pares verdadeiros entre subgrupos, e não de viés intrínseco do classificador na atribuição de escores.
Quando um subgrupo concentra maior prevalência de correspondências reais (como ocorre no quartil Q4 de concordância de nome ou na origem tuberculosa), o VPP eleva-se mecanicamente, mesmo que a TVP permaneça homogênea.
A quase ausência de rejeições na TVP sustenta essa leitura, indicando que o classificador detecta pares verdadeiros com probabilidade similar independentemente do subgrupo, porém a proporção de falsos positivos varia conforme a composição do estrato avaliado.

\section{Discussão}\label{sec:e5-discussao}

\subsection{Equidade do classificador: implicações para a vigilância e limites da análise}\label{sec:e5-discussao-equidade}

A análise de equidade (Seção~\ref{sec:e5-resultados-equidade}) revelou padrão consistente: a taxa de verdadeiros positivos do classificador não variou de forma estatisticamente significativa entre a maioria dos subgrupos avaliados, ao passo que o valor preditivo positivo apresentou disparidades sistemáticas em quatro das cinco dimensões testadas.
A dissociação entre TVP homogênea e VPP heterogêneo possui interpretação direta no contexto de \textit{linkage} probabilístico com prevalência baixa de correspondências verdadeiras \cite{Chouldechova2017fair}: quando a prevalência de pares verdadeiros varia entre subgrupos, o VPP flutua mesmo que a sensibilidade do classificador permaneça constante, porque a proporção de falsos positivos no denominador domina o cálculo em estratos de baixa prevalência.

Essa constatação tem implicação prática para a vigilância epidemiológica da tuberculose. A concentração de verdadeiros positivos no quartil Q4 de concordância de nome e na origem tuberculosa reflete a estrutura subjacente dos dados, na qual pares com alta similaridade de nome e registros provenientes do sistema de notificação de tuberculose apresentam probabilidade intrinsecamente maior de corresponder ao mesmo indivíduo.
O classificador, ao reproduzir essa estrutura, não apresentou indicativos de viés adicional nas dimensões avaliadas e captura a associação genuína entre concordância campo a campo e status de correspondência \cite{Fellegi1969theory,Christen2012book}.

A única rejeição envolvendo TVP (faixa etária 18--39 \textit{versus} 40--59, com $p_{\text{FDR}} = 0{,}039$) situa-se no limiar de significância e pode decorrer de diferenças na qualidade de preenchimento dos campos de identificação entre faixas etárias, hipótese que não pôde ser testada diretamente com os dados disponíveis.
A disparidade de VPP entre sexos (feminino: $0{,}049$; masculino: $0{,}026$; RC~$= 1{,}92$) merece atenção em estudos futuros, particularmente na investigação de possíveis diferenças na completude dos campos de nome entre homens e mulheres nos sistemas de informação em saúde.

A análise também apresenta limitações. A dimensão UF produziu grupo único elegível (RJ), impedindo comparações geográficas que seriam relevantes em país de dimensões continentais.
A faixa etária 0--17 anos, embora com VPP excepcionalmente alto ($0{,}800$), contou com apenas cinco predições positivas, tornando a estimativa instável.
A ausência de variáveis de raça/cor e de geolocalização fina no conjunto de dados utilizado restringe a cobertura da análise de equidade a dimensões disponíveis nos subescores de similaridade, excluindo eixos de estratificação reconhecidamente associados a disparidades em saúde no Brasil \cite{Mehrabi2021fairness}.
A correção de Benjamini-Hochberg \cite{Benjamini1995fdr} controlou a taxa de falsas descobertas no conjunto dos 36 testes, porém a multiplicidade de dimensões e a heterogeneidade de tamanhos amostrais entre grupos exigem cautela na generalização dos achados para bases de dados com composição demográfica distinta.

A análise de equidade encerra os estudos empíricos. A síntese dos achados por objetivo, as contribuições, as limitações e as perspectivas são apresentadas no Capítulo~\ref{cap:conclusao}.

\end{chapter}
